% !TeX spellcheck = en_US
\documentclass[french]{yLectureNote}

\title{Électromagnétisme 2 PS}
\subtitle{Physique}
\author{Paulhenry Saux}
\date{\today}
\yLanguage{Français}
%lionels.calmels@cemes.fr
\professor{M.Brut}
\usepackage{graphicx}%----pour mettre des images
\usepackage[utf8]{inputenc}%---encodage
\usepackage{geometry}%---pour modifier les tailles et mettre a4paper
%\usepackage{awesomebox}%---pour les boites d'exercices, de pbq et de croquis ---d\'esactiv\'e pour les TP de PC
\usepackage{tikz}%---pour deiffner + d\'ependance de chemfig
% \usepackage{tabularx}%---pour dimensionner automatiquement les tableaux avec variable X
\usepackage{awesomebox}%---Pour les boites info, danger et autres
\usepackage{menukeys}%---Pour deiffner les touches de Calculatrice
\usepackage{fancyhdr}%---pour les en-t\^ete personnalis\'ees
\usepackage{blindtext}%---pour les liens
\usepackage{hyperref}%---pour les liens (\`a mettre en dernier)
\usepackage{caption}%---pour la francisation de la l\'egende table vers Tableau
\usepackage{pifont}
\usepackage{array}%---pour les tableaux
\usepackage{lipsum}
\usepackage{yFlatTable}
\usepackage{multicol}
\usepackage{tabularx}
\usepackage{booktabs}
\newcommand{\Lim}[1]{\lim\limits_{\substack{#1}}\:}
\renewcommand{\vec}{\overrightarrow}
\newcommand{\N}[0]{\mathbb{N}}
\newcommand{\dd}{\mathrm{d}}
\newcommand{\norm}[1]{||\vec{#1}||}
\newcommand{\fo}{\psi(\vec{r},t)}
\newcommand{\foe}{\psi(\vec{r},t)\*}
\newcommand{\HH}{\hat{H}}
\newcommand{\hb}{\hbar}
\newcommand{\lap}{\nabla^2}
\newcommand{\lapcc}{\frac{\partial^2 }{\partial x^2}+\frac{\partial^2 }{\partial y^2}+\frac{\partial^2 }{\partial z^2}}
\newcommand{\E}{\vec{E}(M)}
\newcommand{\B}{\vec{B}(M)}
\newcommand{\V}{V(M)}
\newcommand{\er}{\vec{e_{\rho}}}
\newcommand{\ep}{\vec{e_{\varphi}}}
\newcommand{\pip}{\pi^+}
\newcommand{\pim}{\pi^-}
\newcommand{\mv}{\mathcal{V}}
\DeclareMathOperator{\Div}{Div}
\newcommand{\rot}{\vec{\mathrm{rot}}}
\newcommand{\grad}{\vec{\mathrm{grad}}}
\setcounter{chapter}{4}
\begin{document}

\chapter{Ondes EM dans le vide}

\criticalInfo{Rappels : Équations de Maxwell dans le vide}{
En l'absence de sources ($\rho = 0, \vec{j} = \vec{0}$), les équations locales sont :
\begin{itemize}
    \item \textbf{M-G :} $\text{div } \vec{E} = 0$
    \item \textbf{M-T :} $\text{div } \vec{B} = 0$
    \item \textbf{M-F :} $\rot \vec{E} = -\frac{\partial \vec{B}}{\partial t}$
    \item \textbf{M-A :} $\rot \vec{B} = \frac{1}{c^2} \frac{\partial \vec{E}}{\partial t}$ (avec $c = 1/\sqrt{\epsilon_0 \mu_0}$)
\end{itemize}
}

\section{Propagation et Potentiels}

\warningInfo{Équation de d'Alembert}{
Chaque composante des champs $\vec{E}$ et $\vec{B}$ ainsi que des potentiels $\vec{A}$ et $V$ satisfait l'équation d'onde :
$$\lap \vec{A} - \frac{1}{c^2} \frac{\partial^2 \vec{A}}{\partial t^2} = \vec{0}$$
}

\checkInfo{Jauge de Lorenz}{
Pour assurer la propagation des potentiels à la célérité $c$, on impose la condition :
$$\text{div } \vec{A} + \frac{1}{c^2} \frac{\partial V}{\partial t} = 0$$
}

\section{Cinématique des Ondes}

% \begin{definition}[Grandeurs fondamentales]
% On distingue deux périodicités pour une onde progressive monochromatique :
% \begin{itemize}
%     \item \textbf{Période temporelle $T$ (s) :} durée minimale au bout de laquelle le signal se reproduit identiquement en un point donné de l'espace.
%     \item \textbf{Période spatiale $\lambda$ (m) :} distance minimale séparant deux points de l'espace vibrant en phase à un instant donné. On l'appelle \textbf{longueur d'onde}.
% \end{itemize}
% \end{definition}

\checkInfo{Pulsation et Vecteur d'onde}{
\begin{itemize}
    \item \textbf{Pulsation $\omega$ (rad.s$^{-1}$) :} $\omega = \frac{2\pi}{T} = 2\pi f$ (où $f = 1/T$ est la fréquence en Hz).
    \item \textbf{Nombre d'onde $k$ (rad.m$^{-1}$) :} $k = \frac{2\pi}{\lambda}$ (norme du vecteur d'onde $\vec{k}$).
\end{itemize}
}

\criticalInfo{Relations de passage}{
La célérité $c$ fait le lien entre l'espace et le temps :
\begin{equation}
    c = \frac{\lambda}{T} = \lambda f = \frac{\omega}{k}
\end{equation}
Dans le vide, pour une onde électromagnétique, $c = \frac{1}{\sqrt{\epsilon_0 \mu_0}} \approx 3 \cdot 10^8 \text{ m.s}^{-1}$.
}

% \begin{proposition}
% \textbf{Vitesse de phase :} $v_{\phi} = \frac{\omega}{k}$ représente la vitesse de déplacement des surfaces d'onde (phase constante). Dans le vide, elle est égale à $c$ et indépendante de la fréquence (milieu non dispersif).
% \end{proposition}

\tipsInfo{Astuces de calcul}{
En notation complexe $\underline{s} = S_0 e^{i(\omega t - \vec{k} \cdot \vec{r})}$, les opérateurs différentiels se simplifient grandement pour une OPPM :
\begin{itemize}
    \item $\frac{\partial}{\partial t} \longrightarrow i\omega$
    \item $\vec{\nabla} \longrightarrow -i\vec{k}$ (donc $\text{div} \longrightarrow -i\vec{k} \cdot$ et $\vec{\mathrm{rot}} \longrightarrow -i\vec{k} \wedge$)
    \item $\lap \longrightarrow -k^2$
\end{itemize}
}

\section{Ondes Planes Progressives Monochromatiques (OPPM)}

\begin{definition}{OPPM}
Onde dont les surfaces d'onde sont des plans et dont la dépendance temporelle est sinusoïdale de pulsation $\omega$. En notation complexe :
$$\vec{\underline{E}}(\vec{r}, t) = \vec{\underline{E}_0} e^{i(\omega t - \vec{k} \cdot \vec{r})}$$
\end{definition}

\criticalInfo{L'essentiel sur la structure}{
Pour une OPPM se propageant selon $\vec{u}$ ($\vec{k} = k\vec{u}$) :
\begin{itemize}
    \item \textbf{Transversalité :} $\vec{E}$ et $\vec{B}$ sont orthogonaux à $\vec{k}$.
    \item \textbf{Trièdre direct :} $(\vec{u}, \vec{E}, \vec{B})$ forment un trièdre direct.
    \item \textbf{Relation d'amplitude :} $E = cB$ ou $\vec{\underline{B}} = \frac{\vec{u}}{c} \wedge \vec{\underline{E}}$.
    \item \textbf{Phase :} Les deux champs sont en phase.
\end{itemize}
}

\section{Polarisation}

La polarisation est l'étude de la trajectoire du vecteur $\vec{E}$ dans un plan orthogonal à la propagation. Soit $\Delta\phi = \phi_y - \phi_x$ le déphasage.

\begin{table}[h]
\centering
\begin{tabularx}{\textwidth}{@{}l X X@{}}
\toprule
\textbf{Type} & \textbf{Condition sur déphasage $\Delta\phi$} & \textbf{Condition sur amplitudes} \\
\midrule
\textbf{Rectiligne} & $0$ ou $\pi \pmod{2\pi}$ & Quelconques \\
\textbf{Circulaire} & $\pm \pi/2 \pmod{2\pi}$ & $E_{0x} = E_{0y}$ \\
\textbf{Elliptique} & Cas général & Quelconques \\
\bottomrule
\end{tabularx}
\end{table}



\section{Énergie du champ EM}

\begin{proposition}
\textbf{Densité volumique d'énergie :}
$$u_{em} = \epsilon_0 E^2 \quad \text{(équirépartition entre } u_e \text{ et } u_m \text{ pour une OPPM)}$$
\textbf{Vecteur de Poynting :}
$$\vec{\Pi} = \frac{\vec{E} \wedge \vec{B}}{\mu_0} = u_{em} c \vec{u}$$
\end{proposition}

\checkInfo{Moyennes temporelles}{
Pour une OPPM d'amplitude $E_0$ :
$$\langle u_{em} \rangle = \frac{1}{2} \epsilon_0 E_0^2 \quad \text{et} \quad \langle \vec{\Pi} \rangle = \frac{1}{2} \epsilon_0 c E_0^2 \vec{u}$$
}

\section{Ondes stationnaires}

\tipsInfo{Superposition}{
Résulte de la somme de deux ondes identiques se propageant en sens inverse.
\begin{itemize}
    \item \textbf{Découplage :} $\vec{E}(z,t) = 2 E_0 \cos(kz) \cos(\omega t) \vec{e_x}$.
    \item \textbf{N{\oe}uds de E :} $\cos(kz) = 0$.
    \item \textbf{Flux d'énergie :} $\langle \vec{\Pi} \rangle = \vec{0}$ (pas de transport d'énergie).
\end{itemize}
}

\end{document}
